% Vietnamese translation for OI Public Library
% Source-File: IOI中国国家候选队论文/2006/王栋/王栋.doc
\documentclass[11pt]{article}
\usepackage{oipl-vn}
\usepackage{tikz}
\usetikzlibrary{calc}

\newcommand{\Oh}{\mathcal{O}}
\newcommand{\Vor}{\operatorname{Vor}}
\newcommand{\CH}{\operatorname{CH}}
\newcommand{\DT}{\operatorname{DT}}

\title{Voronoi}
\author{Wang Dong}
\date{Đội tuyển quốc gia Trung Quốc, 2006}

\begin{document}
\renewcommand{\figurename}{Hình}
\maketitle

\origtitle{Voronoi}
\sourcepath{IOI China candidate team papers / 2006 / Wang Dong / Wang Dong.doc}

\renewcommand{\abstractname}{Tóm tắt}
\begin{abstract}
Tệp DOC nguồn bị hỏng phần đầu và không còn đọc được như một tài liệu Word
chuẩn. Các đoạn văn, slide đi kèm và mã Pascal trong thư mục nguồn cho thấy
nội dung chính là giới thiệu biểu đồ Voronoi, quan hệ với tam giác hóa
Delaunay, và hai bài ứng dụng \textit{Run Away} cùng \textit{Fat Man}. Bản
dịch này dựng lại phần có thể xác nhận được từ các mảnh văn bản và chương
trình đi kèm, giữ cùng tinh thần trình bày thuật toán hình học tính toán.
\end{abstract}

\noindent\textbf{Từ khóa:} biểu đồ Voronoi; tam giác hóa Delaunay; hình học
tính toán; điểm an toàn nhất.

\section{Khái niệm Voronoi}

Cho tập điểm
\[
  S=\{P_1,P_2,\ldots,P_n\}
\]
trên mặt phẳng. Với hai điểm \(P_i,P_j\), đường trung trực của đoạn
\(P_iP_j\) chia mặt phẳng thành hai nửa mặt phẳng. Gọi
\[
  H(P_i,P_j)=\{X\mid |XP_i|<|XP_j|\}
\]
là nửa mặt phẳng gồm các điểm gần \(P_i\) hơn \(P_j\).

\begin{figure}[ht]
\centering
\begin{tikzpicture}[scale=0.95, every node/.style={font=\small}]
  \fill[blue!7] (0.4,0.2) -- (2.35,4.2) -- (3.25,0.2) -- cycle;
  \draw[gray!55, dashed] (2.35,4.2) -- (3.25,0.2);
  \draw[black!45] (1.1,1.5) -- (4.5,2.25);
  \filldraw[black] (1.1,1.5) circle (1.7pt) node[below left] {\(P_1\)};
  \filldraw[black] (4.5,2.25) circle (1.7pt) node[below right] {\(P_2\)};
  \filldraw[black] (2.8,1.875) circle (1.2pt) node[below] {\(M\)};
  \node[blue!55!black] at (1.35,3.2) {\(H(P_1,P_2)\)};
  \node[gray!55!black, rotate=-77] at (2.55,3.25) {trung trực};
  \draw[->, blue!55!black] (1.65,3.0) -- (2.15,2.55);
\end{tikzpicture}
\caption{Nửa mặt phẳng gần \(P_1\) hơn \(P_2\) được xác định bởi đường trung trực của \(P_1P_2\).}
\end{figure}

Ô Voronoi của \(P_i\) là giao của tất cả các nửa mặt phẳng tương ứng:
\[
  V(P_i)=\bigcap_{j\ne i} H(P_i,P_j).
\]
Nói cách khác, \(V(P_i)\) gồm mọi điểm có điểm gần nhất trong \(S\) là
\(P_i\). Tập các ô \(V(P_i)\) tạo thành biểu đồ Voronoi của \(S\), ký hiệu
\(\Vor(S)\).

Các cạnh của biểu đồ Voronoi nằm trên những đường trung trực. Một đỉnh
Voronoi thường là tâm đường tròn đi qua ba điểm của \(S\); tại đó ba ô
Voronoi gặp nhau. Nếu có bốn điểm đồng viên, biểu đồ có thể suy biến và cần
xử lý cẩn thận trong cài đặt.

\section{Quan hệ với Delaunay}

Tam giác hóa Delaunay \(\DT(S)\) là đồ thị đối ngẫu của biểu đồ Voronoi:
hai điểm \(P_i,P_j\) được nối bởi một cạnh Delaunay khi hai ô Voronoi
\(V(P_i)\) và \(V(P_j)\) kề nhau. Tính chất hình học thường dùng là: một tam
giác thuộc tam giác hóa Delaunay nếu đường tròn ngoại tiếp của nó không chứa
điểm nào khác của \(S\) ở bên trong.

\begin{figure}[ht]
\centering
\begin{tikzpicture}[scale=0.82, every node/.style={font=\small}]
  \clip (-0.1,-0.1) rectangle (7.6,4.6);
  \draw[gray!35] (0,0) rectangle (7.5,4.5);
  \draw[blue!55, thick] (1.7,0) -- (2.0,1.35) -- (2.55,2.1) -- (3.1,4.5);
  \draw[blue!55, thick] (2.0,1.35) -- (0,2.0);
  \draw[blue!55, thick] (2.55,2.1) -- (4.25,2.35) -- (5.65,4.5);
  \draw[blue!55, thick] (4.25,2.35) -- (4.9,0);
  \draw[blue!55, thick] (4.25,2.35) -- (6.9,1.45) -- (7.5,1.7);
  \draw[orange!75!black, semithick, dashed] (1.0,3.45) -- (2.85,2.8) -- (4.0,3.55)
    -- (6.0,3.0) -- (5.7,0.95) -- (3.35,0.85) -- (2.85,2.8);
  \draw[orange!75!black, semithick, dashed] (2.85,2.8) -- (3.35,0.85) -- (1.0,3.45);
  \draw[orange!75!black, semithick, dashed] (4.0,3.55) -- (5.7,0.95) -- (6.0,3.0);
  \filldraw[black] (1.0,3.45) circle (1.7pt) node[above left] {\(P_1\)};
  \filldraw[black] (2.85,2.8) circle (1.7pt) node[above] {\(P_2\)};
  \filldraw[black] (4.0,3.55) circle (1.7pt) node[above] {\(P_3\)};
  \filldraw[black] (6.0,3.0) circle (1.7pt) node[above right] {\(P_4\)};
  \filldraw[black] (3.35,0.85) circle (1.7pt) node[below] {\(P_5\)};
  \filldraw[black] (5.7,0.95) circle (1.7pt) node[right] {\(P_6\)};
  \node[blue!55!black, fill=white, inner sep=1pt] at (1.0,0.55) {cạnh Voronoi};
  \node[orange!75!black, fill=white, inner sep=1pt] at (5.6,0.35) {cạnh Delaunay};
\end{tikzpicture}
\caption{Các cạnh Voronoi và các cạnh Delaunay là hai cấu trúc đối ngẫu trên cùng tập điểm.}
\end{figure}

Quan hệ đối ngẫu này làm Voronoi hữu ích trong nhiều bài toán:
\begin{itemize}
  \item tìm điểm xa nhất hoặc gần nhất theo một tập điểm chướng ngại;
  \item dựng các vùng ảnh hưởng;
  \item xử lý các bài toán khoảng cách cực trị trên mặt phẳng;
  \item chuyển giữa bài toán ô Voronoi và bài toán cạnh Delaunay.
\end{itemize}

\section{Dựng biểu đồ}

Các mảnh trong slide cho thấy hai cách tiếp cận chính: dựng trực tiếp bằng
nửa mặt phẳng và dựng thông qua tam giác hóa Delaunay. Chương trình
\texttt{voronoi.dpr} đi kèm sắp xếp điểm theo tọa độ, dùng các đường trung
trực, tìm giao điểm đường thẳng, và duy trì cạnh kề giữa các điểm. Trong mã có
các thủ tục như:
\begin{itemize}
  \item tìm đường trung trực của hai điểm;
  \item tìm giao của hai đường thẳng;
  \item xác định phía của một điểm đối với một đoạn có hướng;
  \item tìm tiếp tuyến khi chia để trị tập điểm đã sắp xếp.
\end{itemize}

Một khung dựng chia để trị có thể mô tả như sau:
\begin{enumerate}
  \item sắp xếp các điểm theo hoành độ, rồi tung độ;
  \item chia tập điểm thành hai nửa;
  \item dựng biểu đồ Voronoi cho từng nửa;
  \item tìm chuỗi cạnh phân cách hai biểu đồ bằng các tiếp tuyến giữa hai bao
  lồi;
  \item ghép hai biểu đồ và loại các phần cạnh không còn hợp lệ.
\end{enumerate}

Nếu cài đặt đầy đủ và xử lý suy biến tốt, độ phức tạp mục tiêu là
\(\Oh(n\log n)\). Các slide cũng nhắc tới một cách thô hơn có thể lên tới
\(\Oh(n^4)\), thường chỉ phù hợp để minh họa hoặc kiểm chứng dữ liệu nhỏ.

\section{Ví dụ 1: Run Away}

Đề \textit{Run Away} trong mảnh DOC mô tả một căn phòng lớn có nhiều lỗ nhỏ
trên sàn; khi bẫy kích hoạt, ta muốn đứng tại vị trí an toàn nhất, tức vị trí
có khoảng cách tới lỗ gần nhất là lớn nhất. Các dòng còn đọc được cho thấy ví
dụ đầu ra dạng:
\[
  \text{The safest point is }(50.0,50.0).
\]
và một ví dụ khác gần \((1433.0,1669.8)\).

Nếu miền cho phép là một hình chữ nhật và các lỗ là tập điểm \(S\), hàm cần
tối đa hóa là
\[
  d(X)=\min_{P_i\in S}|XP_i|.
\]
Trong mỗi ô Voronoi, điểm gần nhất không đổi, nên cực đại toàn cục không cần
xét ở mọi điểm của mặt phẳng. Ứng viên nằm trong các vị trí sau:
\begin{itemize}
  \item đỉnh Voronoi nằm trong miền;
  \item giao của cạnh Voronoi với biên hình chữ nhật;
  \item các góc của hình chữ nhật.
\end{itemize}

Do đó thuật toán là dựng \(\Vor(S)\), liệt kê các ứng viên hợp lệ, tính
khoảng cách tới điểm gần nhất và chọn giá trị lớn nhất. Một nhận xét trong
slide là nghiệm liên quan tới \(\CH(S)\), bao lồi của các điểm: cạnh không bị
chặn của biểu đồ Voronoi xuất hiện tại các điểm thuộc bao lồi, còn miền hình
chữ nhật giúp cắt các cạnh vô hạn thành đoạn hữu hạn cần kiểm tra.

\begin{figure}[ht]
\centering
\begin{tikzpicture}[scale=0.78, every node/.style={font=\small}]
  \draw[thick] (0,0) rectangle (8,5);
  \filldraw[black] (1.0,1.0) circle (1.6pt) node[below left] {\(P_1\)};
  \filldraw[black] (2.0,4.0) circle (1.6pt) node[above left] {\(P_2\)};
  \filldraw[black] (4.1,2.5) circle (1.6pt) node[below] {\(P_3\)};
  \filldraw[black] (6.4,3.7) circle (1.6pt) node[above right] {\(P_4\)};
  \filldraw[black] (6.8,1.0) circle (1.6pt) node[below right] {\(P_5\)};
  \draw[blue!55, thick] (0,2.2) -- (2.55,2.55) -- (3.6,4.95);
  \draw[blue!55, thick] (2.55,2.55) -- (3.55,1.1) -- (3.1,0);
  \draw[blue!55, thick] (2.55,2.55) -- (5.05,2.7) -- (8,2.35);
  \draw[blue!55, thick] (5.05,2.7) -- (5.4,5);
  \draw[blue!55, thick] (5.05,2.7) -- (5.7,1.45) -- (8,0.65);
  \draw[red!70!black, very thick] (3.05,3.35) circle (2.05);
  \filldraw[red!70!black] (3.05,3.35) circle (1.7pt) node[above right] {ứng viên};
  \filldraw[red!70!black] (5.05,2.7) circle (1.5pt);
  \filldraw[red!70!black] (0,2.2) circle (1.5pt);
  \node[align=center, blue!55!black, fill=white, inner sep=1pt] at (6.35,4.45) {cạnh Voronoi\\cắt biên};
\end{tikzpicture}
\caption{Trong bài \emph{Run Away}, điểm an toàn nhất nằm trong số các đỉnh Voronoi, giao với biên, hoặc góc hình chữ nhật.}
\end{figure}

\section{Ví dụ 2: Fat Man}

Mảnh DOC còn giữ được tên bài \textit{Fat Man}, nguồn Central Europe 1999 và
UVA 295. Ví dụ đầu ra đọc được có dạng:
\[
  \text{Maximum size in test case 1 is }2.2361.
\]
Tệp \texttt{zju2009.dpr} đi kèm thêm bốn góc hình chữ nhật vào tập điểm và
dựng cấu trúc Voronoi. Điều này phù hợp với mô hình: trong một miền hình chữ
nhật có các chướng ngại điểm hoặc lỗ, cần tìm bán kính lớn nhất của một người
hoặc vật thể có thể đi qua/đứng an toàn mà không va vào chướng ngại và biên.

Cách dùng Voronoi giống ví dụ trước. Khoảng cách tới chướng ngại gần nhất chỉ
thay đổi khi vượt qua cạnh Voronoi. Vì vậy các điểm cực trị đáng xét không
nằm tùy ý trong miền mà nằm ở các đỉnh Voronoi, giao với biên, hoặc tại các
điểm biên đặc biệt. Nếu phải xét khoảng cách tới biên hình chữ nhật, có thể
đưa thêm các ràng buộc biên vào phép kiểm tra ứng viên.

\section{Ghi chú cài đặt}

Các chương trình Pascal trong thư mục nguồn cho thấy một hiện thực hình học
tương đối trực tiếp:
\begin{itemize}
  \item dùng \texttt{double} và một hằng \(\varepsilon=10^{-7}\) để so sánh;
  \item lưu điểm với tọa độ và chỉ số ban đầu;
  \item tính đường trung trực bằng trung điểm và một điểm thứ hai trên đường;
  \item biểu diễn đường thẳng dưới dạng \(ax+by+c=0\);
  \item tìm giao điểm bằng giải hệ hai phương trình tuyến tính;
  \item dùng tích có hướng để xác định phía và tìm tiếp tuyến.
\end{itemize}

Khi dịch sang một ngôn ngữ hiện đại, cần chú ý các trường hợp suy biến:
điểm trùng nhau, ba điểm thẳng hàng, bốn điểm đồng viên, cạnh Voronoi song
song, và sai số khi so sánh số thực. Với bài thi, nhiều lỗi không nằm ở ý
tưởng Voronoi mà nằm ở các biên hình học này.

\section{Kết luận}

Biểu đồ Voronoi là một công cụ mạnh để xử lý các bài toán ``gần nhất'' và
``xa nhất'' trên mặt phẳng. Nó chia mặt phẳng thành các vùng mà mỗi vùng có
cùng điểm gần nhất, từ đó biến việc tìm cực trị liên tục thành kiểm tra một
số hữu hạn ứng viên. Qua các ví dụ \textit{Run Away} và \textit{Fat Man}, ta
thấy một mẫu quen thuộc: dựng Voronoi cho các điểm nguy hiểm, cắt với miền
hợp lệ, rồi chọn ứng viên có khoảng cách tới điểm gần nhất lớn nhất.

\section*{Ghi chú về nguồn}

Tệp DOC gốc bắt đầu bằng dữ liệu hỏng và không có chữ ký tệp Word hợp lệ.
Các phần còn đọc được chỉ gồm mảnh văn bản, tên hình, trường nhúng và tên
chương trình. Bản dịch này dựa trên các mảnh còn đọc được trong DOC, slide
\texttt{Wang Dong.ppt}, và hai chương trình \texttt{voronoi.dpr},
\texttt{zju2009.dpr}. Các hình WMF nhúng trong tài liệu gốc không được trích
xuất trực tiếp, nhưng những sơ đồ Voronoi cốt lõi đã được vẽ lại bằng TikZ theo
nội dung còn xác nhận được từ slide và văn bản.

\end{document}
